\documentclass[11pt,a4paper]{article}

\usepackage[margin=1in]{geometry}
\usepackage[T1]{fontenc}
\usepackage[utf8]{inputenc}
\usepackage{lmodern}
\usepackage{microtype}
\usepackage{booktabs}
\usepackage{array}
\usepackage{longtable}
\usepackage{graphicx}
\usepackage{float}
\usepackage{hyperref}
\usepackage{xcolor}

\newcommand{\placeholderfigure}[1]{%
  \fbox{\parbox[c][4.2cm][c]{0.92\linewidth}{\centering \textit{#1}}}%
}

\newcommand{\metric}[1]{\texttt{\detokenize{#1}}}
\newcommand{\metriclines}[1]{\begin{tabular}[t]{@{}l@{}}#1\end{tabular}}

\title{Deep Learning for Football Match Analysis with YOLOv8 and Multi-Object Tracking}
\author{Alessandro ARENSBERG}
\date{\today}

\begin{document}

\maketitle

\begin{abstract}
This report presents a technically structured football video analysis pipeline built on YOLOv8, ByteTrack-style tracking, colour-based team assignment, camera-motion compensation, and ball recovery heuristics. The goal is not only to detect objects, but to convert broadcast footage into stable trajectories and match-level statistics. On the evaluated 750-frame clip, the fine-tuned pipeline detects the ball in 715 frames (95.3\%), compared with 670 frames (89.3\%) for a generic baseline, while reducing the total number of player tracks (from 190 to 46) and detecting referees consistently for the fine-tuned runs. An ablation study shows that ball interpolation is the dominant factor for coverage, adding 51 frames over the no-interpolation variant.
\end{abstract}

\tableofcontents
\newpage

\section{Introduction}
Football broadcast analysis is a challenging video understanding problem because the relevant objects are small, fast, partially occluded, and observed through a moving camera. The project goal is to transform a single broadcast clip into stable object trajectories, team assignments, ball possession estimates, and readable visual analytics.

The technical core of the system is a representation-learning pipeline built on YOLOv8. The detector is fine-tuned on football-specific annotations and then combined with temporal post-processing: ByteTrack-style association, jersey-colour clustering, ball interpolation, ROI recovery, camera-motion compensation, and approximate pitch geometry. The report is intentionally structured around this pipeline so that the design choices, metrics, and limitations are explicit.

\subsection{Contributions}
\begin{itemize}
  \item A modular football analysis pipeline that separates detection, tracking, team inference, recovery, geometry, and visualisation.
  \item A comparison between a generic YOLOv8 baseline and a football fine-tuned detector.
  \item A focused ablation study that isolates the effect of ball interpolation and track memory.
  \item A failure analysis that identifies what remains difficult at the detector level and what requires better calibration or data.
\end{itemize}

\section{Problem Setting}
Given a broadcast football video, the system must detect four classes (ball, player, goalkeeper, referee), maintain identities over time, infer team identity for players, and convert frame-level observations into match-level quantities. The evaluation is carried out on the same 750-frame clip for all variants so that differences are attributable to the pipeline design rather than to the input data.

This setting is especially difficult for the ball, which is both tiny and motion-blurred, and for referees, whose visual appearance overlaps with players. In addition, camera motion introduces apparent displacement that must be compensated before any metric computed in image space can be interpreted geometrically.

\section{Dataset and Annotation Protocol}
The dataset is a football broadcast collection exported in the standard YOLO format. The target classes are \metric{ball}, \metric{player}, \metric{goalkeeper}, and \metric{referee}. The release is split into train, validation, and test subsets, and the fine-tuning stage uses the training split while the held-out subsets are reserved for model selection and qualitative checks.

\begin{table}[H]
\centering
\caption{Dataset summary used for detector fine-tuning.}
\begin{tabular}{lccc}
\toprule
Split & Images / frames & Labels & Role \\
\midrule
Train & 612 & 612 & Fine-tuning \\
Validation & 38 & 38 & Model selection \\
Test & 13 & 13 & Sanity check \\
\bottomrule
\end{tabular}
\end{table}

The annotation distribution is highly imbalanced: players dominate the frames, referees are much rarer, and the ball is the most difficult class because of its size and frequent occlusion. This imbalance motivates both transfer learning and temporal recovery.

\section{Methodology}

\subsection{Detector Architecture and Transfer Learning}
The detector is a YOLOv8 model whose convolutional backbone learns hierarchical visual representations: low-level edges and textures in early layers, more semantic football-specific patterns in deeper layers. Fine-tuning adapts these representations to broadcast football footage, where the key cues differ from generic image datasets.

The baseline configuration uses generic pretrained YOLOv8 weights with minimal post-processing. The full configuration uses football fine-tuned weights and the complete downstream pipeline.

In practice, fine-tuning improves the detector in three ways. First, it adapts to the pitch viewpoint and broadcast framing. Second, it specialises the model for small objects such as the ball. Third, it reduces confusion between similar classes, especially player versus referee.

\subsection{Multi-Object Tracking}
Detections are associated over time using a ByteTrack-style tracker. The role of tracking is not only to assign IDs, but also to stabilise trajectories, reduce fragmentation, and enable track-level statistics such as track length and identity consistency.

To improve robustness, the implementation keeps a short track memory. If a detection disappears briefly and reappears with sufficient overlap, the previous identity can be recovered instead of starting a new track. This is a simple but effective temporal prior.

\subsection{Team Assignment and Ball Possession}
Team identity is inferred from jersey-colour descriptors extracted from player crops. The system clusters colour features and propagates team labels along tracks. This design is intentionally lightweight: it avoids training another classifier while still providing usable team-level analytics.

Ball possession is estimated geometrically by comparing the ball position to the lower part of nearby player boxes. The logic is heuristic, but it is appropriate for this project because the ball is often missing for short intervals and the available signal is spatial rather than semantic.

\subsection{Ball Interpolation and ROI Recovery}
The ball is the hardest object to track consistently. Short gaps are filled by linear interpolation, and when the detector misses the ball, a region-of-interest recovery pass revisits the last known area. These mechanisms do not invent new evidence; they simply exploit temporal continuity to bridge short occlusions.

\subsection{Camera Motion and Pitch Geometry}
Camera motion compensation is estimated with optical flow on stable regions of the frame and used as a frame-to-frame translation correction. The corrected positions are then mapped to pitch coordinates through an approximate perspective transform. This step is necessary because any distance estimate in image coordinates is distorted by camera motion and perspective.

\subsection{Metrics}
The evaluation focuses on metrics that reflect both detection quality and downstream stability. Ball coverage measures the fraction of frames in which the ball is observed or interpolated. Track fragmentation counts how often identities are broken into multiple segments. Referee detection is tracked separately because it is a useful proxy for domain adaptation. Team-switch counts measure whether player identities remain stable enough for aggregation.

\begin{table}[H]
\centering
\caption{Metric families used in the analysis.}
\resizebox{\textwidth}{!}{%
\begin{tabular}{p{0.30\linewidth}p{0.60\linewidth}}
\toprule
Metric family & Interpretation \\
\midrule
Detection coverage & Frames with ball, player, goalkeeper, referee detections \\
Tracking stability & Total track count, track fragmentation, memory reactivation \\
Recovery quality & Observed versus interpolated ball frames, maximum gap length \\
Team consistency & Average and maximum team switches per track \\
Geometry quality & Camera offset stability and approximate field-space distances \\
\bottomrule
\end{tabular}%
}
\end{table}

\section{Experimental Protocol}
All variants are evaluated on the same broadcast clip. The goal is to isolate the effect of each module while keeping the input video, sampling rate, and post-processing thresholds fixed.

\begin{table}[H]
\centering
\caption{Configurations compared in the experiments.}
\resizebox{\textwidth}{!}{%
\begin{tabular}{llcccc}
\toprule
Variant & Weights & Camera motion & Interpolation & Team assignment & ROI recovery \\
\midrule
baseline & generic YOLOv8 & off & off & off & off \\
finetuned\_full & fine-tuned football & on & on & on & on \\
finetuned\_no\_interp & fine-tuned football & on & off & on & on \\
finetuned\_no\_cam & fine-tuned football & off & on & on & on \\
finetuned\_no\_team & fine-tuned football & on & on & off & on \\
\bottomrule
\end{tabular}%
}
\end{table}

\section{Results}

\subsection{Main Comparison}
The main comparison highlights an important point: raw detection count is not enough. The generic baseline detects the ball in slightly more frames, but it fragments player identities much more heavily and fails completely on the referee class.

\begin{table}[H]
\centering
\caption{Baseline versus full pipeline on the 750-frame evaluation clip.}
\resizebox{\textwidth}{!}{%
\begin{tabular}{lrrrr}
\toprule
Variant & Ball frames & Player tracks (total) & Referee frames & Qualitative stability \\
\midrule
baseline & 670/750 (89.3\%) & 190 & 0/750 (0\%) & Fragmented \\
finetuned\_full & 715/750 (95.3\%) & 46 & 750/750 (100\%) & Stable \\
\bottomrule
\end{tabular}%
}
\end{table}

The fine-tuned model is preferable because the downstream tracking stage becomes much more coherent. The improvement is especially visible for referees, where transfer learning adapts the detector to a class that the generic model misses entirely.

\subsection{Ablation Study}
The ablation study shows that ball interpolation is the strongest contributor to ball coverage. Track memory also helps, but its effect is smaller than interpolation.

\begin{table}[H]
\centering
\caption{Ablation study for ball coverage and temporal recovery.}
\resizebox{\textwidth}{!}{%
\begin{tabular}{lrr}
\toprule
Variant & Ball frames & Loss relative to full pipeline \\
\midrule
finetuned\_full & 715/750 (95.3\%) & -- \\
finetuned\_no\_interp & 664/750 (88.5\%) & -51 frames (-6.8\%) \\
finetuned\_no\_cam & 715/750 (95.3\%) & 0 frames on coverage, worse geometry \\
finetuned\_no\_team & 715/750 (95.3\%) & 0 frames on coverage, weaker aggregation \\
\bottomrule
\end{tabular}%
}
\end{table}

\subsection{Qualitative Results}
The system produces an annotated output video in which detections, tracks, team colours, formation lines, and minimap overlays remain readable even when the camera pans. The qualitative result is important because it reveals whether the pipeline is visually stable, not just numerically strong.

\begin{figure}[H]
\centering
\includegraphics[width=\textwidth]{screenshot-baseline.png}
\caption{Baseline output placeholder.}
\end{figure}

\begin{figure}[H]
\centering
\includegraphics[width=\textwidth]{screenshot-with-predictions.png}
\caption{Fine-tuned full pipeline output example.}
\end{figure}

\section{Deep Learning Analysis}
The report explicitly links the system to deep learning concepts from the course. YOLOv8 is a convolutional network whose backbone learns transferable feature hierarchies. Fine-tuning on the football dataset changes the internal representations so that they become more sensitive to the appearance of football-specific objects, especially the ball and referees.

The results show that representation learning matters more than raw model size in this setting. The generic detector appears strong on the ball in frame-level coverage, but its features are not specialised enough to support stable tracking or referee recognition. By contrast, the fine-tuned model produces fewer fragmented trajectories and a much more useful class distribution for downstream analysis.

This also explains why the temporal modules matter. A detector alone predicts boxes independently per frame, whereas the project needs a coherent sequence representation. Interpolation, track memory, and team propagation turn per-frame CNN outputs into usable temporal signals.

\section{Failure Analysis}
The hardest failures are not random; they correspond to specific modelling limitations.

\begin{table}[H]
\centering
\caption{Failure modes and technical causes.}
\resizebox{\textwidth}{!}{%
\begin{tabular}{p{0.23\linewidth}p{0.30\linewidth}p{0.37\linewidth}}
\toprule
Failure mode & Observed symptom & Main cause \\
\midrule
Goalkeeper ambiguity & Switches between player and goalkeeper labels & Class similarity and insufficient specialised data \\
Ball blur / occlusion & Short gaps and missed detections & Small object scale, motion blur, and occlusion \\
Camera panning & Apparent jitter in positions and overlays & Frame-to-frame motion not fully compensated \\
Calibration error & Speed and distance estimates are approximate & No exact homography / pitch reference points \\
\bottomrule
\end{tabular}%
}
\end{table}

The goalkeeper problem is the clearest example of a limit that cannot be solved with post-processing alone. The detector needs more specialised training data, or a better class definition, to separate players from goalkeepers consistently. Similarly, ball recovery is limited by the absence of physics: a frame-based detector can bridge gaps, but it cannot fully infer a fast airborne trajectory.

\section{Discussion and Limitations}
The main takeaway is that football analysis is a systems problem. Detection quality matters, but so do temporal smoothing, camera compensation, and geometric consistency. The best results come from composing several simple modules rather than trusting a single end-to-end prediction.

The current implementation also has limits. First, the evaluation is based on one broadcast clip, so the generalisation gap is not fully measured. Second, metric conversion depends on an approximate pitch model, which makes speed estimates less reliable than detection metrics. Third, the system does not yet include an automatic line-based calibration procedure, which would improve field-space accuracy.

Future work should therefore focus on three directions: more football-specific data for the detector, better calibration for metric estimation, and a more principled temporal model for the ball. These changes would improve both accuracy and robustness.

\section{Conclusion}
This project delivers a structured football analysis pipeline that connects YOLOv8 detection, multi-object tracking, team inference, and temporal recovery into a coherent system. The fine-tuned model improves downstream stability, detects referees reliably, and produces cleaner trajectories than a generic baseline. More importantly, the ablation and failure analysis show why each module matters and where the remaining technical bottlenecks lie.

\begin{thebibliography}{9}
\bibitem{youtube_tutorial} Public tutorial that inspired the project: \textit{Build an AI/ML Football Analysis system with YOLO, OpenCV, and Python}
\newline \url{https://www.youtube.com/watch?v=neBZ6huolkg}
\bibitem{yolov8} Ultralytics. YOLOv8 Documentation.\newline \url{https://docs.ultralytics.com/models/yolov8} -- \url{https://yolov8.com/}
\bibitem{bytetrack} Y. Zhang et al. \textit{ByteTrack: Multi-Object Tracking by Associating Every Detection Box}\newline \url{https://arxiv.org/abs/2110.06864}
\bibitem{football} \textit{football-players-detection} dataset, Roboflow export.\newline \url{https://universe.roboflow.com/roboflow-jvuqo/football-players-detection-3zvbc/dataset/1}
\end{thebibliography}

\end{document}